


\title{Constructions in Quasi-Categories}

\maketitle





\begin{tikzcd}
& 3 \\
& 0 \ar[dl] \ar[dr] \ar[u] \\
1 \ar[uur, bend left] \ar[rr] & & 2 \ar[uul, bend right]
\end{tikzcd}



\begin{tikzcd}
& 0 \ar[dl, swap, "f"] \ar[dr, "fg"] \\
1 \ar[rr, swap, "g"] & & 2 
\end{tikzcd}










\begin{tikzcd}
& 3 \\
& 0 \ar[dl, swap, "f"] \ar[dr, "fg"] \ar[u, "fgh"] \\
1 \ar[uur, bend left, "gh"] \ar[rr, swap, "g"] & & 2 \ar[uul, bend right, swap, "h"]
\end{tikzcd}





\begin{definition}{Homotopy}{•}
$f \simeq g$ if there is a $2$-morphism of the form
\begin{center}
\begin{tikzcd}
& 0 \ar[dl, swap, "id"] \ar[dr, "g"] \\
1 \ar[rr, swap, "f"] & & 2 
\end{tikzcd}

\end{center}
\end{definition}





\begin{lemma}{•}{•}

\end{lemma}
\begin{proof-lemma}



\newthought{Reflexive: }


\begin{tikzcd}
& 0 \ar[dl, swap, "id"] \ar[dr, "f"] \\
1 \ar[rr, swap, "f"] & & 2 
\end{tikzcd}


\newthought{Symmetric: } 

%
%\begin{tikzcd}[column sep=scriptsize]
%A \arrow[dr] \arrow[rr, "f"{name=U, below, draw=red}]{}
%& & B \arrow[dl] \\
%& C \ar[from=U]
%\end{tikzcd}
%%


%\begin{tikzcd}
%A \arrow[to=Z, red] \arrow[to=2-2, blue]
%& B \\
%|[alias=Z]| C
%& D
%\arrow[from=ul, to=1-2]
%\end{tikzcd}


\begin{tikzcd}
& 3 \\
& 0 \ar[dl, swap,"id"] \ar[dr, "id"] \ar[u, "f"] \\
1 \ar[uur, bend left, "g"{name=U}] \ar[rr, swap, "id"]{} \ar[from=U, to=2-2, phantom, "\blacksquare"] & & 2 \ar[uul, bend right, swap, "f"] 
\end{tikzcd}



\newthought{Transitive: } 

\begin{tikzcd}
& 3 \\
& 0 \ar[dl, swap, "id"] \ar[dr, "id"] \ar[u, "h"] \\
1 \ar[uur, bend left, "g"] \ar[rr, swap, "id"] & & 2 \ar[uul, bend right, swap, "f"{name=U}] \ar[from=U, to=2-2, phantom, "\blacksquare"]
\end{tikzcd}



\end{proof-lemma}








\begin{definition}{Composition}{•}


any choice of $fg$

\begin{tikzcd}
& 0 \ar[dl, swap, "f"] \ar[dr, "fg"] \\
1 \ar[rr, swap, "g"] & & 2 
\end{tikzcd}

\end{definition}




\section{Fibrations}




\begin{definition}{Fibre of a Functor}{•}

\begin{tikzcd}
\cat{F} \ar[r, "\mid" marking] \ar[d, "\mid" marking] & \cat{B} \ar[d] \\
* \ar[r] & \cat{C}
\end{tikzcd}
\end{definition}
\begin{defframe}

since simplicial sets are presheafs, their pullbacks are done pointwise

So we just need pullbacks in sets

so given a functor, which maps $B_i \to C_i$, we just apply the pullback on that

So the fibre is the elements of $B_i$ that get mapped to the identity


\end{defframe}





\begin{prop}{•}{•}

is an acyclic kan fibration

so each of the fibers are contractible

(fibers are the pullback
\end{prop}



\begin{lemma}{•}{•}

\end{lemma}
\begin{proof-lemma}



\begin{tikzcd}
& 3 \\
& 0 \ar[dl, swap,"id"] \ar[dr, "id"] \ar[u, "f"] \\
1 \ar[uur, bend left, "g"{name=U}] \ar[rr, swap, "id"]{} \ar[from=U, to=2-2, phantom, "\blacksquare"] & & 2 \ar[uul, bend right, swap, "f"] 
\end{tikzcd}



\end{proof-lemma}



\begin{definition}{Homotopy Category}{•}

\end{definition}




\begin{definition}{Mapping Space}{•}

\end{definition}


\begin{prop}{•}{•}
is a kan complex
\end{prop}
\begin{proof-prop}



\begin{tikzcd}
& 3 \\
& 0 \ar[dl, swap,"id"] \ar[dr, "id"] \ar[u, "h"] \\
1 \ar[uur, bend left, "g"{name=U}] \ar[rr, swap, "id"]{} \ar[from=U, to=2-2, phantom, "\blacksquare"] & & 2 \ar[uul, bend right, swap, "f"] 
\end{tikzcd}




\end{proof-prop}




\begin{definition}{$\infty$-Groupoid}{•}
all morphisms are equivalences
\end{definition}


\begin{prop}{•}{•}
iff kan complex
\end{prop}
\begin{proof-prop}


from 
we know that the mapping space is a kan complex, so we can fill every horn 



\end{proof-prop}







\begin{definition}{Functor}{•}

natural transformation
\end{definition}











